
\section{Topología computacional, cronología parcial e identificabilidad}

\subsection{Espacio de estados y complejos duales}

Sea \(D\subset\mathbb R^3\) el dominio espacial de interés. Para un estado inverso \(\tau\), definimos el dominio material
\[
M_\tau\subset D
\]
y su complemento libre
\[
V_\tau=D\setminus M_\tau.
\]

Una discretización computacional puede representarse mediante complejos cubicales o simpliciales
\[
K_\tau^+,\qquad K_\tau^-,
\]
asociados respectivamente a materia y espacio negativo. Esta construcción no pretende introducir una nueva noción de homología; la contribución del marco consiste en usar ambos complejos de forma acoplada para restringir historias constructivas.

\subsection{Homología relativa respecto del exterior}

Una dificultad inmediata aparece cuando cámaras y corredores permanecen conectados al exterior. La conectividad ordinaria puede entonces ocultar identidad arquitectónica.

Sea
\[
E_\tau\subset V_\tau
\]
la componente o subregión exterior distinguida. Consideramos la pareja
\[
(V_\tau,E_\tau)
\]
y sus grupos de homología relativa
\[
\boxed{
H_k(V_\tau,E_\tau;\mathbb F).
}
\]

El objetivo es describir estructura interna respecto del exterior, manteniendo explícitas las condiciones de frontera.

\subsection{Filtración por despeje}

Definimos
\[
d_M(x,\tau)
=
\operatorname{dist}(x,M_\tau)
\]
y, para un radio \(r\ge0\),
\[
\boxed{
V_{\tau,r}
=
\{
x\in V_\tau:
d_M(x,\tau)\ge r
\}.
}
\]

Para \(\tau\) fijo y \(r_1\le r_2\),
\[
V_{\tau,r_2}\subseteq V_{\tau,r_1},
\]
de modo que \(r\) sí induce una filtración ordinaria. Esto permite estudiar cuándo conexiones estrechas desaparecen mientras recintos de mayor escala permanecen topológicamente distinguibles \cite{edelsbrunner2010,deywang2022}.

Los números de Betti
\[
\beta_k(\tau,r)
=
\dim H_k(V_{\tau,r};\mathbb F)
\]
resumen componentes, ciclos y cavidades.

\subsection{Estabilidad frente a perturbaciones geométricas}

Para funciones tame adecuadas \(f,g\), la estabilidad de diagramas de persistencia establece
\[
\boxed{
d_B(D(f),D(g))
\le
\|f-g\|_\infty,
}
\]
bajo las hipótesis del teorema de estabilidad \cite{cohensteiner2007}. Esta propiedad resulta especialmente relevante para reconstrucciones provenientes de escáneres, nubes de puntos y volúmenes discretizados con ruido.

\subsection{Corrección formal: la coordenada inversa no genera necesariamente una bifiltración}

Una familia \((\tau,r)\mapsto V_{\tau,r}\) sólo constituye una bifiltración clásica cuando existen inclusiones monotónicas en ambos parámetros.

Durante construcción o desmontaje, un espacio puede abrirse, cerrarse, fusionarse con el exterior o separarse. Por tanto, en general,
\[
V_{\tau_i,r}
\nsubseteq
V_{\tau_{i+1},r}
\]
y tampoco necesariamente ocurre la inclusión inversa.

En consecuencia, IRL adopta una formulación zigzag para la dimensión temporal:
\[
\boxed{
V_{0,r}
\leftrightarrow
V_{1,r}
\leftrightarrow
\cdots
\leftrightarrow
V_{n,r}.
}
\]

Aplicando homología:
\[
H_k(V_{0,r})
\leftrightarrow
H_k(V_{1,r})
\leftrightarrow
\cdots
\leftrightarrow
H_k(V_{n,r}),
\]
lo cual se inscribe en el marco de zigzag persistence \cite{carlssondesilva2010}. Esta precisión corrige una formulación anterior del proyecto: \((\tau,r)\mapsto V_{\tau,r}\) se denomina ahora \emph{familia de dos parámetros} salvo que se demuestre monotonicidad suficiente para hablar de bifiltración.

\subsection{Persistencia multiparámetro}

Cuando se utilizan dos parámetros monotónicos, por ejemplo despeje \(r\) y umbral de densidad \(\rho_c\), puede definirse
\[
\mathcal V_{r,\rho_c}
=
\{
x:
d_M(x)\ge r,\;
\rho(x)\le\rho_c
\}.
\]

Bajo una orientación coherente de inclusiones,
\[
(r,\rho_c)
\mapsto
H_k(\mathcal V_{r,\rho_c})
\]
forma un módulo de persistencia multiparámetro. Este régimen conecta naturalmente geometría y muografía, aunque sus invariantes son más complejos que un barcode unidimensional \cite{carlssonzomorodian2009}.

\subsection{Grafos de Reeb y transiciones críticas}

Para una función
\[
f:X\rightarrow\mathbb R,
\]
el grafo de Reeb identifica puntos que pertenecen a una misma componente conexa de un conjunto de nivel. Puede utilizarse con altura, distancia al exterior, despeje, densidad o posterior espacial de vacío \cite{biasotti2008,deywang2022}.

En un modelo suave, candidatos a transición topológica aparecen en puntos críticos:
\[
\nabla f(x^\star)=0.
\]

En una formulación implícita:
\[
\boxed{
\Phi(x^\star;\lambda^\star)=0,
\qquad
\nabla_x\Phi(x^\star;\lambda^\star)=0.
}
\]

IRL utiliza esta estructura como detector de estados donde cambia la clase de accesibilidad.

\subsection{Orden parcial y politopo de cronologías}

Sea
\[
P=(E,\prec)
\]
el poset constructivo. Una historia serial compatible es una extensión lineal
\[
\pi\in\mathcal L(P).
\]

Asignando a cada evento un tiempo normalizado
\[
\mathbf t=(t_1,\ldots,t_n)\in[0,1]^n,
\]
se define el politopo de orden de Stanley
\[
\boxed{
\mathcal O(P)
=
\{
\mathbf t\in[0,1]^n:
t_i\le t_j
\text{ si }e_i\prec e_j
\}.
}
\]

Para un poset finito de \(n\) elementos,
\[
\boxed{
\operatorname{Vol}(\mathcal O(P))
=
\frac{e(P)}{n!},
}
\]
donde \(e(P)=|\mathcal L(P)|\) \cite{stanley1986}. Esto proporciona una medida geométrica del grado en que las precedencias reducen el espacio de cronologías normalizadas.

Definimos como métrica descriptiva del modelo:
\[
C_P
=
1-\frac{e(P)}{n!}.
\]

\(C_P\) cuantifica compresión causal dentro del poset elegido y no debe interpretarse como probabilidad histórica.

\subsection{Accesibilidad en espacio de configuraciones}

Para una pieza rígida con configuración
\[
q=(R,p)\in SE(3),
\]
el espacio libre es
\[
\mathcal C_{\mathrm{free}}(t)
=
\{
q:
B(q)\cap M(t)=\varnothing
\}.
\]

La colocación exige una curva continua
\[
\gamma:[0,1]\rightarrow\mathcal C_{\mathrm{free}}(t)
\]
y, en el marco IRL, restricciones adicionales:
\[
\gamma(s)
\in
\mathcal C_{\mathrm{free}}
\cap
\mathcal C_{\mathrm{force}}
\cap
\mathcal C_{\mathrm{support}}.
\]

La abstracción de espacio de configuraciones es estándar en motion planning \cite{lavalle2006}; su función aquí es transformar geometría de acceso en restricciones cronológicas constructivas.

\subsection{Reducción del espacio de historias}

Sea \(\mathcal H_0\) la familia inicial de historias y
\[
\mathcal H_j
=
\{
H\in\mathcal H_{j-1}:g_j(H)=1
\}
\]
el conjunto superviviente después de un gate.

Para una medida \(\mu\),
\[
\boxed{
R_j
=
1-
\frac{\mu(\mathcal H_j)}
{\mu(\mathcal H_{j-1})}.
}
\]

\(\mu\) puede ser conteo, volumen del politopo, masa posterior o medida Monte Carlo. Esta cantidad convierte una pregunta cualitativa en un objetivo experimental: medir cuánto restringe cada modalidad el espacio de historias.

\subsection{Entropía posterior y métrica descriptiva de identificabilidad}

Para historias discretas \(H_i\) con probabilidades posteriores \(p_i\),
\[
\mathsf H(H\mid Y)
=
-\sum_i p_i\log p_i.
\]

Como resumen descriptivo se propone
\[
\boxed{
I_{\mathrm{IRL}}
=
1-
\frac{\mathsf H(H\mid Y)}
{\log n}.
}
\]

Esta cantidad se reporta junto con la distribución completa: es una métrica de comunicación del proyecto, no un sustituto de análisis formal de identificabilidad.

\subsection{Diseño experimental óptimo}

Para un diseño de medición \(d\), la ganancia esperada de información es
\[
\boxed{
\operatorname{EIG}(d)
=
\mathbb E_{Y_d}
\left[
D_{\mathrm{KL}}
\left(
p(H,\Theta\mid Y,Y_d,d)
\|
p(H,\Theta\mid Y)
\right)
\right].
}
\]

La siguiente medición puede elegirse mediante
\[
d^\star
=
\arg\max_d
\operatorname{EIG}(d),
\]
sujeto a costo, seguridad y viabilidad \cite{alexanderian2021}. En Khufu, este formalismo puede comparar nuevas posiciones muográficas, escaneo geométrico, caracterización litológica o geofísica subsuperficial.

\subsection{Implicación para el dashboard}

La implementación v1.2 distingue tres niveles:

\begin{enumerate}
\item grafo de conectividad como aproximación explicativa;
\item familia de despeje como modelo morfológico;
\item homología volumétrica/persistente como objetivo numérico siguiente.
\end{enumerate}

Los valores \(\beta_0^{\mathrm{graph}}\) y
\[
\beta_1^{\mathrm{graph}}
=
|E|-|V|+\beta_0^{\mathrm{graph}}
\]
mostrados en la interfaz describen el grafo, no la homología 3D completa del monumento.
